%
% reversi_180_lbfgs.tex
%
% This file is part of the reversi program
% http://github.com/rcrr/reversi-v2
%
% Author Roberto Corradini mailto:rob_corradini@yahoo.it
% Copyright 2026 Roberto Corradini. All rights reserved.
%
% License
%
% This program is free software; you can redistribute it and/or modify it
% under the terms of the GNU General Public License as published by the
% Free Software Foundation; either version 3, or (at your option) any
% later version.
%
% This program is distributed in the hope that it will be useful,
% but WITHOUT ANY WARRANTY; without even the implied warranty of
% MERCHANTABILITY or FITNESS FOR A PARTICULAR PURPOSE. See the
% GNU General Public License for more details.
%
% You should have received a copy of the GNU General Public License
% along with this program; if not, write to the Free Software
% Foundation, Inc., 59 Temple Place - Suite 330, Boston, MA  02111-1307, USA
% or visit the site <http://www.gnu.org/licenses/>.
%
%
% Capitolo 10 — L'algoritmo L-BFGS: il solutore numerico del livello OPTIMIZE,
% trattato in pieno: approssimazione inversa dell'Hessiano a memoria limitata,
% ricorsione a doppio ciclo, ricerca del passo con condizioni di Wolfe forti
% e interpolazione di Hermite cubica.
%

\chapter{L'algoritmo L-BFGS}

\label{ch:lbfgs}

Il capitolo 9 ha definito la funzione di perdita $J(\bm{w})$ — l'errore
quadratico medio con regolarizzazione di ridge — e il suo gradiente, e ha
spiegato \emph{perché} la minimizzazione non può essere affidata al metodo di
Newton: l'Hessiana di $J$ avrebbe dimensioni incompatibili con la memoria
disponibile (centinaia di migliaia di pesi per i pattern più ricchi), mentre
il gradiente si calcola a ogni iterazione senza problemi. Il capitolo 6 ha
quindi introdotto il livello \texttt{OPTIMIZE} della pipeline, che esegue
l'ottimizzazione dei pesi con un algoritmo a \emph{memoria limitata} della
famiglia BFGS — in sigla L-BFGS — senza però svilupparne la trattazione
matematica: tutto quel dettaglio è stato qui rimandato, a questo capitolo.

Questo è dunque il capitolo di riferimento per il solutore: contiene la
trattazione matematica \emph{completa} dell'algoritmo implementato nel modulo
\texttt{reversi/tools/opt\_lbfgs.py} — l'approssimazione inversa dell'Hessiana
costruita implicitamente sulla finestra dei $m$ passi più recenti
(\/\emph{two-loop recursion}), la ricerca del passo a ogni iterazione
(condizioni di Wolfe \emph{forti}, espansione guidata da interpolazione cubica
di Hermite, fase di \emph{zoom}) e la descrizione puntuale delle scelte di
implementazione, ciascuna collegata alla formula che essa realizza.

\section{Perché la memoria limitata}

Sia $\bm{x}$ il vettore dei parametri da stimare (nel nostro caso il vettore
di pesi $\bm{w}$) e $f(\bm{x})$ il funzionale da minimizzare, di classe
$C^1$. L'idea dei metodi \emph{quasi-Newton} è sostituire la matrice
$[\nabla^2 f(\bm{x}_k)]^{-1}$ del metodo di Newton con una matrice
$\bm{H}_k$ mantenuta ricorsivamente: a ogni passo si usa la direzione
\begin{equation}
\label{eq:quasi-newton-direction}
\bm{p}_k = -\,\bm{H}_k\,\bm{g}_k,
\qquad
\bm{g}_k = \nabla f(\bm{x}_k),
\end{equation}
e si aggiorna $\bm{H}_k$ a $k{+}1$ usando solo le differenze di gradiente
osservate. Il metodo di BFGS classico fa esattamente questo, ma mantiene
\emph{esplicita} una matrice $n\times n$: $O(n^2)$ memoria per matrice e
$O(n^2)$ operazioni per moltiplicazione $\bm{H}_k\bm{g}_k$. Con $n$ nell'ordine
di $10^5$--$10^6$ l'approccio è fuori portata.

L'osservazione di cui si alimenta la variante a memoria limitata (Liu,
Storey, 1989; Byrd, Nocedal, 1995) è che l'aggiornamento BFGS ha una
struttura di \emph{somma di rank due} che può essere applicata a un vettore
\emph{senza mai formare la matrice}: il prodotto $\bm{H}_{k+1}\bm{g}$ si
calcola a partire dal prodotto $\bm{H}_k\bm{g}$, da un passo $\bm{s}$ e da una
differenza di gradiente $\bm{y}$, con costi $O(n)$ per vettore coinvolto. Se
invece di aggiornare una matrice completa si tengono solo gli $m$ \emph{coppie}
più recenti $(\bm{s}_i,\bm{y}_i)$ — con $m \ll n$, tipicamente $5 \le m \le 20$
— si ottiene un'approssimazione $\bm{H}_k$ di rango effettivo $2m$, rappresentata
da $2m$ vettori ($O(mn)$ memoria) e applicabile a un gradiente con $O(mn)$
operazioni. Per $m=50$ e $n \sim 10^5$ si passa da $\sim 10^{10}$ coefficienti
(matrice piena) a $\sim 10^{7}$ (la finestra): questo è il motivo per cui
l'ottimizzazione dei pesi della pipeline è affidata a L-BFGS.

\section{L'aggiornamento BFGS e la sua forma implicita}

Partiamo dall'aggiornamento che L-BFGS rappresenta implicitamente. Al passo
$k$ l'algoritmo compie la mossa
\begin{equation}
\bm{x}_{k+1} = \bm{x}_k + \alpha_k \bm{p}_k,
\end{equation}
dove $\alpha_k > 0$ è il passo scelto dalla ricerca del passo (\/\S\ref{sec:lbfgs-line-search})
e definisce le due grandezze fondamentali:
\begin{equation}
\label{eq:s-y}
\bm{s}_k = \bm{x}_{k+1} - \bm{x}_k,
\qquad
\bm{y}_k = \bm{g}_{k+1} - \bm{g}_k,
\qquad
\bm{g}_{k+1} = \nabla f(\bm{x}_{k+1}).
\end{equation}
La coppia $(\bm{s}_k,\bm{y}_k)$ soddisfa la \emph{condizione secante}:
\begin{equation}
\label{eq:secant}
\bm{H}_{k+1}\,\bm{y}_k = \bm{s}_k,
\end{equation}
cioè l'approssimazione $\bm{H}_{k+1}$ dell'inversa dell'Hessiana deve
``riprodurre'' il passo effettivamente compiuto quando la applica alla
differenza di gradiente. Tra tutte le matrici che soddisfano la (\ref{eq:secant}),
l'aggiornamento BFGS è quella che minimizza la distanza (in una norma pesata
dall'Hessiana precedente) da $\bm{H}_k$; la sua forma canonica è
\begin{equation}
\label{eq:bfgs-update}
\bm{H}_{k+1} =
(\bm{I} - \rho_k \bm{s}_k \bm{y}_k^{\!\top})\,
\bm{H}_k\,
(\bm{I} - \rho_k \bm{y}_k \bm{s}_k^{\!\top})
\;+\; \rho_k\, \bm{s}_k \bm{s}_k^{\!\top},
\qquad
\rho_k = \frac{1}{\bm{y}_k^{\!\top}\bm{s}_k}.
\end{equation}
Due dettagli che l'implementazione rispetta:

\begin{itemize}
\item l'esistenza di $\rho_k$ richiede $\bm{y}_k^{\!\top}\bm{s}_k > 0$;
  questa è la \emph{condizione di curvatura}, e vale in ogni vicinato
  dell'ottimo per una funzione convessa. Il codice impone la condizione con
  una soglia esplicita (\/\S\ref{sec:lbfgs-history}): una coppia che non la
  soddisfa viene \emph{scartata}, non archiviata;
\item nell'implementazione il fattore di scala di $\bm{H}_0$ non è l'identità
  ma $\gamma \bm{I}$ con $\gamma$ ricalcolato a ogni iterazione
  (\/\S\ref{sec:lbfgs-twoloop}): è una variante precisa di cui va resa
  esplicita la formula.
\end{itemize}

\section{La ricorsione a doppio ciclo}
\label{sec:lbfgs-twoloop}

L'apporto centrale della forma a memoria limitata è questo: se $\bm{H}_0$ è
sufficientemente semplice, il prodotto $\bm{H}_k\bm{g}$ — e quindi la
direzione di ricerca (\ref{eq:quasi-newton-direction}) — si calcola
ricorsivamente sulle $m$ coppie più recenti, senza mai costruire
$\bm{H}_k$. Iterando $m$ volte la (\ref{eq:bfgs-update}) dal punto $k-m$ si
ottiene, per $i = k-1, k-2, \dots, k-m$ indicizzate nella finestra:
\begin{equation}
\label{eq:twoloop}
\bm{H}_k = \bm{V}_m\,\bigl(\gamma \bm{I}\bigr)\,\bm{V}_m^{\!\top},
\end{equation}
dove $\bm{V}_m$ è il prodotto in catena degli $m$ fattori a forma di
$(\bm{I}-\rho_i\bm{s}_i\bm{y}_i^{\!\top})(\bm{I}-\rho_i\bm{y}_i\bm{s}_i^{\!\top})
+\rho_i\bm{s}_i\bm{s}_i^{\!\top}$ di (\ref{eq:bfgs-update}). Svolgendo il
prodotto e applicandolo a $\bm{g}$ si ottiene la \emph{ricorsione a doppio
ciclo} (two-loop recursion), che è esattamente ciò che implementa la funzione
\texttt{\_two\_loop\_recursion} di \texttt{opt\_lbfgs.py}:

\begin{enumerate}
\item \textbf{primo ciclo} (dall'ultimo verso il primo): si parte da
  $\delta \leftarrow \bm{g}$ e, per $i = m-1, \dots, 0$,
  \begin{equation}
  \label{eq:twoloop-loop1}
  \alpha_i = \rho_i\, \bm{s}_i^{\!\top}\delta,
  \qquad
  \delta \leftarrow \delta - \alpha_i\, \bm{y}_i;
  \end{equation}
\item \textbf{scalatura}: si applica il fattore centrale
  \begin{equation}
  \label{eq:twoloop-gamma}
  \bm{r} \leftarrow \gamma\,\delta,
  \qquad
  \gamma =
  \begin{cases}
  \dfrac{\bm{s}_{m-1}^{\!\top}\bm{y}_{m-1}}
         {\bm{y}_{m-1}^{\!\top}\bm{y}_{m-1}}, & m > 0,\\[2mm]
  1, & m = 0,
  \end{cases}
  \end{equation}
  cioè $\gamma\bm{I}$ è la scelta di $\bm{H}_0$: non l'identità, ma la matrice
  diagonale che \emph{interpoli la secante più recente} (è il classico
  scalare $\gamma$ di L-BFGS, qui ricalcolato a ogni iterazione dalla coppia
  più recente della finestra);
\item \textbf{secondo ciclo} (dal primo verso l'ultimo): per $i = 0, \dots, m-1$,
  \begin{equation}
  \label{eq:twoloop-loop2}
  \beta_i = \rho_i\, \bm{y}_i^{\!\top}\bm{r},
  \qquad
  \bm{r} \leftarrow \bm{r} + (\alpha_i - \beta_i)\, \bm{s}_i.
  \end{equation}
\end{enumerate}

Il vettore $\bm{r}$ risultante è $\bm{H}_k\bm{g}$; la direzione di ricerca
è la sua negata:
\begin{equation}
\label{eq:p-direction}
\bm{p}_k = -\,\bm{H}_k\bm{g}_k.
\end{equation}
Con finestra vuota ($m=0$) la (\ref{eq:twoloop-gamma}) dà $\gamma = 1$ e la
ricorsione degenera in $\bm{H}_k\bm{g} = \bm{g}$: la direzione è la
discesa del gradiente, come deve essere al primo passo.

Il listing mostra il codice reale (sostanzialmente integrale; i commenti
spiegativi del modulo sono rimossi):

\begin{lstlisting}[language=python, caption={La ricorsione a doppio ciclo, da
\texttt{opt\_lbfgs.py}: la struttura dei due cicli e la scalatura $\gamma$.}]
def _two_loop_recursion(g, s, y, rho):
    m = len(s)
    alpha = np.zeros(m)

    prod = np.empty(g.shape, dtype=np.float64)   # buffer prodotto, preallocato

    delta = g.copy()
    for i in reversed(range(m)):
        alpha[i] = rho[i] * np.dot(s[i], delta)  # (5): alpha_i
        np.multiply(y[i], alpha[i], out=prod)
        np.subtract(delta, prod, out=delta)      # (5): delta

    if m > 0:
        last_s = s[-1]; last_y = y[-1]
        gamma = np.dot(last_s, last_y) / np.dot(last_y, last_y)
    else:
        gamma = 1.0
    r = gamma * delta                            # (6)

    for i in range(m):
        beta = rho[i] * np.dot(y[i], r)          # (7)
        np.multiply(s[i], alpha[i] - beta, out=prod)
        np.add(r, prod, out=r)                   # (7): r
    return r
\end{lstlisting}

I costi sono quelli promessi: $2m$ prodotti scalari e $4m$ vettori- scalar
di dimensione $n$, quindi $O(mn)$ operazioni e $O(mn)$ memoria per la
finestra. Il confronto con Newton è il motivo dell'esistenza di tutto il
capitolo: $O(mn)$ contro $O(n^2)$ a passo, e $O(mn)$ contro $O(n^2)$ di
memoria.

\section{La ricerca del passo}
\label{sec:lbfgs-line-search}

Fissa la direzione $\bm{p}_k$ (che deve essere \emph{di discesa},
$\bm{p}_k^{\!\top}\bm{g}_k < 0$; il caso contrario è gestito dal riavvio,
\/\S\ref{sec:lbfgs-restart}), il problema della ricerca del passo è
unidimensionale: minimizzare la funzione restringente
\begin{equation}
\label{eq:phi}
\varphi(\alpha) = f(\bm{x}_k + \alpha\,\bm{p}_k),
\qquad \alpha \ge 0.
\end{equation}
L'implementazione non cerca il minimo esatto (richiederebbe troppe
valutazioni di $f$ e del gradiente) ma un punto che soddisfi le
\emph{condizioni di Wolfe forti}, con le due costanti $c_1, c_2$ della
configurazione (default $c_1 = 10^{-4}$, $c_2 = 0{.}9$, con il vincolo
$0 < c_1 < c_2 < 1$ imposto dallo schema di configurazione):
\begin{equation}
\label{eq:wolfe}
\varphi(\alpha) \le \varphi(0) + c_1\,\alpha\,\varphi'(0)
\qquad\text{(sufficiente discesa, Armijo)}
\end{equation}
\begin{equation}
\label{eq:wolfe-curv}
\bigl|\varphi'(\alpha)\bigr| \le -\,c_2\,\varphi'(0)
\qquad\text{(curvatura)}.
\end{equation}
Qui $\varphi'(0) = \bm{g}_k^{\!\top}\bm{p}_k < 0$ e
$\varphi'(\alpha) = \nabla f(\bm{x}_k + \alpha\bm{p}_k)^{\!\top}\bm{p}_k$.
L'implementazione, la funzione \texttt{\_strong\_wolfe\_line\_search},
segue la struttura classica in \emph{due fasi}: una fase di
\emph{espansione} (si cresce $\alpha$ finché il punto non è candidato) e una
fase di \emph{zoom} (si stringe un intervallo che contiene il candidato).
Entrambe usano come prossimo punto di prova il minimo dell'interpolante
cubica di Hermite costruita sui valori noti — l'interpolazione di cui al
\/\S\ref{sec:lbfgs-hermite} — con un numero massimo di iterazioni fissi.

\subsection{La fase di espansione}

Si parte da $\alpha = 1$ (\emph{non} da un valore adattivo: il passo
unitario è il primo tentativo) e si valuta $\varphi$ e $\varphi'$; i valori
già conosciuti all'origine sono quelli di $\bm{x}_k$, riutilizzati senza
ricalcolo. A ogni iterazione $i$ l'algoritmo verifica, in questo ordine,
(codice di \texttt{\_strong\_wolfe\_line\_search}):

\begin{enumerate}
\item se $\varphi(\alpha_1) > \varphi(0) + c_1\alpha_1\varphi'(0)$
  \emph{oppure}, per $i>0$, se $\varphi(\alpha_1) \ge \varphi(\alpha_{\mathrm{prev}})$
  (la funzione smette di scendere): si è ``superato'' il minimo
  raggiungibile lungo la retta — si passa alla fase di zoom sull'intervallo
  $[0, \alpha_1]$;
\item se $|\varphi'(\alpha_1)| \le -c_2\varphi'(0)$: il punto soddisfa le
  Wolfe forti — \textbf{accettato}, la ricerca si conclude;
\item se $\varphi'(\alpha_1) \ge 0$: la derivata è già tornata positiva
  prima ancora che la curvatura fosse soddisfatta — zoom sull'intervallo
  $[\alpha_1, 0]$ (passato con estremi invertiti, la fase di zoom li
  riordina internamente).
\end{enumerate}

Se nessuna condizione si verifica e il budget di espansione non è esaurito
(default \texttt{max\_iters} $= 3$), si sceglie il prossimo $\alpha$:
\begin{equation}
\label{eq:expansion-next}
\alpha_{\mathrm{new}} =
\min\bigl(\hat{\alpha}_{\mathrm{cub}},\; \alpha_{\max}\bigr),
\qquad \alpha_{\max} = 10,
\end{equation}
dove $\hat{\alpha}_{\mathrm{cub}}$ è il minimo dell'interpolante cubica di
Hermite sull'intervallo $[0, \alpha_1]$ (\/\S\ref{sec:lbfgs-hermite});
se l'interpolante non ha un minimo stazionario valido in intervallo, la
regola di fallback è il \emph{raddoppio}: $\alpha_{\mathrm{new}} = 2\alpha_1$.
A budget esaurito il punto corrente viene accettato comunque: è una
safety-valve che garantisce che la ricerca del passo non diverga mai in
valutazioni.

In sintesi: l'origine è già in cache, la fase di espansione costa al più
3 valutazioni di $\varphi$ (il punto $\alpha = 1$ e al più due punti
interpolati), la fase di zoom al più 20.

\subsection{La fase di zoom}
\label{sec:lbfgs-zoom}

La funzione \texttt{\_zoom} riceve due estremi $a_{\mathrm{lo}}, a_{\mathrm{hi}}$
(il caso di estremi invertiti viene normalizzato all'inizio di ogni
iterazione: $x_0 = \min$, $x_1 = \max$) e un budget di
\texttt{max\_iters} $= 20$. A ogni iterazione:

\begin{enumerate}
\item si leggono (dalla cache, v. oltre) $\varphi$ e $\varphi'$ ai due
  estremi;
\item si calcola il minimo $\hat{\alpha}$ dell'interpolante cubica di
  Hermite su $[x_0, x_1]$; se l'interpolante non fornisce un minimo valido
  (nessuna radice reale, o nessuna con curvatura positiva, o fuori
  intervallo), il punto di prova è il \emph{punto medio}:
  \begin{equation}
  \label{eq:zoom-next}
  \alpha = \hat{\alpha}
  \quad\text{se valido},\qquad
  \alpha = \tfrac{1}{2}(x_0 + x_1)\ \text{altrimenti};
  \end{equation}
\item si valutano $\varphi(\alpha)$ e $\varphi'(\alpha)$ e si applica il
  test di accettazione e la riduzione dell'intervallo:
  \begin{itemize}
  \item se $\varphi(\alpha) > \varphi(0) + c_1\alpha\varphi'(0)$
    \emph{oppure} $\varphi(\alpha) \ge \varphi(a_{\mathrm{lo}})$:
    $a_{\mathrm{hi}} \leftarrow \alpha$;
  \item altrimenti, se $|\varphi'(\alpha)| \le -c_2\varphi'(0)$:
    \textbf{accettato};
  \item altrimenti, se $\varphi'(\alpha)(a_{\mathrm{hi}} - a_{\mathrm{lo}})
    \ge 0$: $a_{\mathrm{hi}} \leftarrow a_{\mathrm{lo}}$ (il minimo è
    per forza oltre $\alpha$ nella direzione opposta); in ogni caso
    $a_{\mathrm{lo}} \leftarrow \alpha$.
  \end{itemize}
\end{enumerate}

A budget esaurito si accetta l'ultimo $\alpha$ valutato (stessa
safety-valve dell'espansione). La struttura è quella della ricerca del
minimo unidimensionale ``cubica- con-ricaduta-sul-mediano'' usata dai
sottoroutine di line search di Brent/Press et al.; le uniche differenze
rispetto a quelle referenze sono la cadenza fissa di accettazione per
sopraffusione di budget e l'assenza di un'interpolazione quadratica
separata: il codice usa \emph{solo} la cubica di Hermite e il punto medio.

\subsection{L'interpolazione cubica di Hermite}
\label{sec:lbfgs-hermite}

Il cuore numerico della ricerca del passo è la funzione
\texttt{\_hermite\_cubic\_interp}, che costruisce e valuta l'interpolante
cubica $h$ che coincide in \emph{valore e derivata prima} con $\varphi$ agli
estremi $x_0, x_1$ di un intervallo. Posto
$h = x_1 - x_0$, $t = (x - x_0)/h$ e denotando
$y_0 = \varphi(x_0)$, $y_1 = \varphi(x_1)$,
$dy_0 = \varphi'(x_0)$, $dy_1 = \varphi'(x_1)$, l'interpolante è
\begin{equation}
\label{eq:hermite}
h(x) = y_0\,h_0(t) + y_1\,h_1(t) + dy_0\,h\,h_2(t) + dy_1\,h\,h_3(t),
\end{equation}
con i quattro \emph{polinomi base di Hermite}
\begin{equation}
\label{eq:hermite-basis}
\begin{aligned}
h_0(t) &= 2t^3 - 3t^2 + 1, & h_1(t) &= -2t^3 + 3t^2,\\
h_2(t) &= t^3 - 2t^2 + t,  & h_3(t) &= t^3 - t^2.
\end{aligned}
\end{equation}
Questi polinomi soddisfano le condizioni di cardinalità
$h_j(i) = \delta_{ji}$ e $h_j'(i) = 0$ per $j = 0,1$ e
$h_2'(0) = h_3'(1) = 1$, $h_2'(1) = h_3'(0) = 0$: per questo
(\ref{eq:hermite}) interpola esattamente valori e derivate ai due nodi.
La funzione di codice valuta inoltre le derivate prima e seconda di $h$
(stesse combinazioni lineari dei polinomi base derivati) perché
serviranno a decidere se un punto stazionario è un minimo.

La funzione \texttt{\_hermite\_cubic\_interp\_min} trova il minimo
dell'interpolante in $[x_0, x_1]$ derivando (\ref{eq:hermite}):
$h'(x) = 0$ è equivalente a un'equazione quadratica in $t$,
\begin{equation}
\label{eq:hermite-quadratic}
a\,t^2 + b\,t + c = 0,
\end{equation}
con i coefficienti calcolati dal codice:
\begin{equation}
\label{eq:hermite-coeffs}
\begin{aligned}
a &= -6\,dy + 3\,h\,(dy_0 + dy_1),\\
b &= \;\;6\,dy - 2\,h\,(2\,dy_0 + dy_1),\\
c &= \;\;\;\;\;\;h\,dy_0,
\end{aligned}
\qquad dy = y_1 - y_0.
\end{equation}
La procedura di selezione è:
\begin{itemize}
\item se $a = 0$ (l'interpolante è al più quadratica in $t$) oppure il
  discriminante $b^2 - 4ac \le 0$: nessun punto stazionario reale —
  funzione non utile, si ritorna all'algoritmo chiamante che cadrà sul
  fallback (raddoppio o punto medio);
\item altrimenti si calcolano le due radici
  $t = \dfrac{-b \pm \sqrt{b^2 - 4ac}}{2a}$, si mappano in
  $x = x_0 + t\,h$ e si \emph{verificano}: solo una radice con
  $h''(x) > 0$ è un minimo locale. Il codice testa la prima e poi la
  seconda; se nessuna ha curvatura positiva, ritorna `Non è.
\end{itemize}
Nota che la verifica `curvatura positiva è essenziale: l'interpolante
cubica può avere un massimo e un minimo in intervallo, e usare il massimo
come prossimo passo di prova porterebbe la ricerca nella direzione
sbagliata.

\subsection{La cache delle valutazioni}

Un dettaglio di implementazione che ha un ruolo matematico preciso: la
funzione $\varphi$ usata dalla ricerca del passo è \emph{chiusa su una
cache} a chiavi $\alpha$:
\begin{lstlisting}[language=python, caption={La cache per-$\alpha$ dentro
\texttt{\_strong\_wolfe\_line\_search}.}]
cache = {}

def phi(alpha):
    if alpha not in cache:
        f, grad = fg(x + alpha * p)
        cache[alpha] = FG(f, grad)
    fg_val = cache[alpha]
    return fg_val.fun, np.dot(fg_val.grad, p)
\end{lstlisting}
Ogni valutazione di $f$ al punto $\bm{x} + \alpha\bm{p}$ costa, in questo
problema, un passaggio completo sulla matrice di disegno (tutte le
posizioni del set di addestramento): è la valutazione più costosa di
tutta la run. La cache garantisce che un $\alpha$ già visto (l'origine
$\alpha = 0$ viene messa in cache \emph{prima} di iniziare, riutilizzando
il gradiente appena calcolato nell'iterazione principale) non venga mai
ricalcolato: la fase di zoom, che rilegge continuamente i due estremi
dell'intervallo, non produce quindi \emph{alcuna} valutazione extra su di
essi. Il costo reale di una ricerca del passo è il numero di $\alpha$
\emph{distinti} valutati: 1 (l'origine, gratuita) $+$ al più 3
(espansione) $+$ al più 20 (zoom).

\section{La gestione della finestra $(s, y)$}
\label{sec:lbfgs-history}

Alla fine di ogni iterazione l'algoritmo aggiorna la finestra di memoria.
Il codice, nel corpo principale di \texttt{lbfgs}, compie:

\begin{lstlisting}[language=python, caption={L'aggiornamento della finestra
di memoria, da \texttt{lbfgs}.}]
x_new = x + alpha * p
s = x_new - x
y = fgv_new.grad - fgv.grad

if np.dot(y, s) > CURVATURE_THRESHOLD:   # 1e-15
    if len(sl) == m:
        sl.pop(0); yl.pop(0); rho.pop(0)
    sl.append(s); yl.append(y)
    rho.append(1.0 / np.dot(y, s))
\end{lstlisting}

Le regole sono tre, e ciascuna corrisponde a un requisito matematico:

\begin{enumerate}
\item \textbf{condizione di curvatura}: la coppia $(\bm{s},\bm{y})$ viene
  archiviata solo se $\bm{y}^{\!\top}\bm{s} > 10^{-15}$ (la costante
  \texttt{CURVATURE\_THRESHOLD}). È la guardiola numerica che assicura
  che $\rho = 1/(\bm{y}^{\!\top}\bm{s})$ sia ben definita e che
  l'aggiornamento (\ref{eq:bfgs-update}) mantenga
  $\bm{H}_{k+1}$ definita positiva; una coppia con
  $\bm{y}^{\!\top}\bm{s} \le 0$ (possibile lontano dall'ottimo, su funzioni
  non convesse) viene \emph{scartata} senza far avanzare la finestra;
\item \textbf{rotazione FIFO}: a finestra piena ($m$ coppie) la coppia più
  \emph{vecchia} viene rimossa (\texttt{pop(0)}) prima di inserire la
  nuova: la finestra contiene sempre le $m$ coppie \emph{più recenti},
  come richiesto dalla formulazione (\ref{eq:twoloop});
\item \textbf{archiviazione di $\rho$}: lo scalare
  $\rho_i = 1/(\bm{y}_i^{\!\top}\bm{s}_i)$ viene calcolato e memorizzato
  \emph{una volta sola}, al momento dell'inserzione, e non ricalcolato a
  ogni uso nella ricorsione: la (\ref{eq:twoloop-loop1}) e la
  (\ref{eq:twoloop-loop2}) lo usano come dato.
\end{enumerate}

Le tre liste \texttt{sl}, \texttt{yl}, \texttt{rho} (una per componente
della coppia) restano sempre \emph{allineate} e della stessa lunghezza
$\le m$; questo invariante è garantito a costruzione e ri-verificato al
salvaggio dei checkpoint (\/\S\ref{sec:lbfgs-checkpoint}).

\section{Il riavvio e i criteri di arresto}
\label{sec:lbfgs-restart}

\subsection{Il riavvio}

Prima della ricerca del passo, l'algoritmo verifica che la direzione
calcolata sia \emph{di discesa}:
\begin{equation}
\bm{p}_k^{\!\top}\bm{g}_k \ge 0
\quad\Longrightarrow\quad
\text{riavvio}.
\end{equation}
In tal caso la finestra viene \emph{vuotata} ($\bm{H}_k$ degrada
all'identità, $\gamma = 1$) e la direzione torna la discesa del gradiente,
$\bm{p}_k = -\bm{g}_k$. È la clausola di \emph{restart} classica di L-BFGS:
l'approssimazione $\bm{H}_k$ non è più definita positiva (ad esempio per
una sequenza di coppie scartate dalla condizione di curvatura, o per
accumulo di errore in float32 lontano dall'ottimo) e proseguire con una
$\bm{H}_k$ malcondizionata sarebbe controproducente; ripartire dalla
discesa del gradiente ricostruisce progressivamente una finestra sana.

\subsection{I due criteri di convergenza}

La run si arresta per uno di tre motivi, codificati nel campo
\texttt{reason} del dizionario di informazione restituito:

\begin{itemize}
\item \texttt{CONVERGED\_MIN\_GRAD}: la norma del gradiente
  $\|\bm{g}_k\| < \epsilon_g$ per $N_g$ iterazioni \emph{consecutive}
  (default $\epsilon_g = 10^{-5}$, $N_g = 3$; la coppia è il parametro
  \texttt{min\_grad} $= [\epsilon_g, N_g]$). È il criterio di
  \emph{stazionarietà} classico: in un minimo locale il gradiente tende a
  zero. Il requisito ``per $N_g$ iterazioni consecutive'' (contatore
  azzerato alla prima iterazione non conforme) rende il criterio robusto
  contro le oscillazioni transitorie della norma;
\item \texttt{CONVERGED\_LOW\_PROGRESS}: il \emph{calo relativo} della
  funzione
  \[
  \frac{f_{k-1} - f_k}{|f_k|} < \epsilon_f
  \]
  per $N_f$ iterazioni consecutive (default
  $\epsilon_f = 10^{-14}$, $N_f = 7$; la coppia è il parametro
  \texttt{min\_p\_fun\_decrease} $= [\epsilon_f, N_f]$; ``p'' sta per
  \emph{percentage} in senso di frazione). È un criterio di
  \emph{stall}: la run non sta più scendendo in modo misurabile, anche se
  il gradiente non è ancora sotto soglia. Particolare: per $f_k = 0$ il
  calo relativo non è definito e il contatore viene \emph{congelato}
  (né incrementato né azzerato), così da non far divergere né
  falsificare il test;
\item \texttt{MAX\_ITERS\_REACHED}: il budget di iterazioni
  \texttt{max\_iters} (default 100 nell'API; 500 nella configurazione di
  pipeline) è esaurito senza che i due criteri sopra si siano verificati.
\end{itemize}

L'ordine di valutazione a ogni iterazione $k$ è: prima il test sul
gradiente (sul punto corrente, \emph{prima} di calcolare la direzione),
poi il test sul calo relativo (che confronta il punto corrente con il
precedente, tramite la variabile \texttt{fgv\_fun\_prev} aggiornata a
fine passo), infine — solo se nessuno dei due ha fermato la run — la
calcolo della direzione, la ricerca del passo e l'aggiornamento.

\section{Scelte di implementazione}
\label{sec:lbfgs-impl}

Ogni scelta del modulo \texttt{opt\_lbfgs.py} è il riflesso di una formula
del \S\S\ref{sec:lbfgs-twoloop}--\ref{sec:lbfgs-hermite}; le elenchiamo
collegandole esplicitamente.

\subsection{La struttura del modulo e il contratto \texttt{FG}}

Il modulo è \emph{autosufficiente}: non importa nulla dalla pipeline
(né \texttt{LMContext}, né database, né livelli). L'unica cosa che
chiede al chiamante è una funzione \texttt{fg} che, dato un vettore
$\bm{x}$, restituisca la coppia (valore di $f$, gradiente). La coppia è
incapsulata nel named-tuple \texttt{FG} $= (\texttt{fun},\texttt{grad})$,
che è il tipo con cui il modulo dialoga internamente (la cache della
ricerca del passo, il riporto del risultato). Questa separazione è ciò
che permette al medesimo solutore di essere usato sia dal livello
\texttt{OPTIMIZE} della pipeline (dove \texttt{fg} è la closure
gradiente del modello logistico) sia dai test unitari su funzioni
sintetiche (quadratiche, Rosenbrock, Zakharov).

\subsection{I tipi di dato: il contratto float32}

\texttt{lbfgs} accetta $\bm{x}_0$ solo in \texttt{float32} o
\texttt{float64} (altrimenti \texttt{TypeError}) e ne fa una
\emph{copia} (\texttt{x = x0.copy()}), così che l'input del chiamante non
venga mai mutato. Il contratto di pipeline che questa scelta serve è
preciso: il vettore di pesi viaggia in \texttt{float32} dall'inizio alla
fine della run, e l'ottimizzatore \emph{deve} restituire un vettore
\texttt{float32} (la verifica è nel test
\texttt{test\_float32\_end\_to\_end}: \texttt{x}, tutti i $\bm{s}$ e
$\bm{y}$ archiviati, e la norma del gradiente restano float32).
Internamente la ricorsione a doppio ciclo opera con una semantica di dtype
particolare, che il sotto-paragrafo successivo spiega.

\subsection{Il buffer di prodotto preallocato (identità bit-per-bit)}

Il listing 1 mostra il buffer \texttt{prod}, allocato \emph{una volta per
chiamata} in \texttt{float64} e riusato a ogni operazione vettoriale
intermedia tramite la forma \texttt{out=} di NumPy:
\texttt{np.multiply(y[i], alpha[i], out=prod)} seguito da
\texttt{np.subtract(delta, prod, out=delta)}. Il motivo è una questione
di \emph{identità numerica}, non solo di performance:

\begin{itemize}
\item il prodotto di uno scalare \texttt{float64} (il $\alpha_i$) per un
  array \texttt{float32} produce in NumPy un array \texttt{float64};
  l'assegnazione in-place su un buffer \texttt{float32}
  (\texttt{delta -= alpha * y}) fa quindi \emph{primiero il prodotto in
  f64 e poi il downcast in f32} — esattamente la semantica che il buffer
  \texttt{prod} in f64 + \texttt{out=} replica;
\item l'allocazione di $2m$ array temporanei \emph{freschi} a ogni
  chiamata (il comportamento ``classico'' di
  \texttt{delta -= alpha * y}) era il vero collo di bottiglia a scala
  reale: con $m = 93$ e $n = 77943$ in f32 si trattava di $\sim 58$ MB
  di allocazioni per chiamata. Il buffer singolo preallocato elimina
  l'overhead di allocazione \emph{senza cambiare un solo bit} del
  risultato: il test \texttt{test\_out\_equals\_original\_bit\_identical}
  fissa l'identità bit-per-bit contro l'implementazione di riferimento,
  su tre configurazioni seed/dimensione/dtype (f32 $n=50$ $m=5$;
  f32 $n=500$ $m=9$; f64 $n=64$ $m=7$).
\end{itemize}

La ricorsione è implementata in \emph{NumPy puro}, senza kernelli
compilati: a scala reale il collo di bottiglia non è il conteggio delle
operazioni ma l'allocazione dei temporanei vettoriali (con $m$ nell'ordine
di una decina di coppie e $n$ di $10^5$ componenti, $2m$ prodotti
vettoriali freschi a ogni chiamata ammontano a decine di MB), e il buffer
preallocato con semantica bit-per-bit è l'intero meccanismo di
performance: la traiettoria L-BFGS non cambia, solo la memoria
temporanea smette di essere riallocata.

\subsection{La configurazione del solutore}

La tabella raccoglie i parametri di \texttt{lbfgs} con il loro ruolo
matematico e il default di pipeline (lo schema pydantic
\texttt{hlp\_config\_schemas.py} li espone nella sezione
\texttt{optimization} del file di configurazione):

\begin{table}[htbp]
\centering
\small
\begin{tabular}{llll}
\toprule
Parametro & Ruolo matematico & Default API & Default pipeline \\
\midrule
\texttt{c1}        & Costante di Armijo, (\ref{eq:wolfe}) & $10^{-4}$ & $10^{-4}$ \\
\texttt{c2}        & Costante di curvatura, (\ref{eq:wolfe-curv}) & $0.9$ & $0.9$ \\
\texttt{min\_grad} & $[\epsilon_g, N_g]$ del test di stazionarietà & $[10^{-5}, 3]$ & dalla config \\
\texttt{min\_p\_fun\_decrease} & $[\epsilon_f, N_f]$ del test di stall & $[10^{-14}, 7]$ & dalla config \\
\texttt{max\_iters} & Budget assoluto di iterazioni $k$ & $100$ & $500$ \\
\texttt{m}         & Lunghezza della finestra di memoria, (\ref{eq:twoloop}) & $10$ & $50$ \\
\texttt{log\_every\_n} & Cadenza di logging (0 = silenzioso) & $0$ & dalla config \\
\texttt{save\_every\_n} & Cadenza di checkpoint (0 = nessuno) & $0$ & dalla config \\
\bottomrule
\end{tabular}
\caption{Parametri di \texttt{lbfgs} con il ruolo matematico di ciascuno.}
\end{table}

Lo schema di configurazione \emph{impone} $c_2 > c_1$ (validazione
pydantic): la disuguaglianza è la condizione di coerenza delle Wolfe
forti, non un capriccio dell'interfaccia.

\subsection{Logging e conteggio delle valutazioni}

Il logging (cadenza \texttt{log\_every\_n}, sink \texttt{logger}
iniettato, default \texttt{print}) riporta a ogni iterazione loggata: il
valore di $f$, la norma del gradiente, il $\alpha$ accettato, l'occupazione
della finestra $\,|\texttt{sl}|/m\,$ e il \emph{numero di chiamate a
\texttt{fg} dall'ultima riga di log}. Il contatore di chiamate
(\texttt{fg\_call\_count}) è incrementato da un wrapper
\texttt{fgc\_plain} che circonda la funzione del chiamante: è la misura
diretta del costo reale della run in valutazioni del funzionale, il
quantità che conta quando si confrontano due configurazioni. La riga
finale riporta il totale.

\subsection{Checkpoint e warm start}
\label{sec:lbfgs-checkpoint}

L'ottimizzazione dei pesi è la fase più lunga della pipeline (ordine di
grandezza: minuti a ore, a seconda della scala del modello); un crash a
metà run non deve far perdere tutto. Per questo \texttt{lbfgs} espone due
meccanismi simmetrici:

\begin{itemize}
\item \textbf{salvaggio periodico}: a ogni
  $k \bmod \texttt{save\_every\_n} = 0$ (con $k > 0$) viene invocato il
  callback \texttt{save\_fn} con la firma a 7 argomenti
  \texttt{save\_fn}$(\bm{x}, \texttt{sl}, \texttt{yl}, \texttt{rho}, k, f,
  g_{\mathrm{norm}}^{\max})$. Il sesto argomento non è la norma del
  gradiente corrente ma il \emph{massimo} di una finestra scorrevole
  (\texttt{FixedSizeMaxList}) delle ultime $N_g$ norme: è la quantità
  che il criterio \texttt{CONVERGED\_MIN\_GRAD} monitora, e salvarla
  ``come lo vede il criterio'' rende il checkpoint auto-consistente.
  Nel livello \texttt{OPTIMIZE} il callback è
  \texttt{hlp\_optimize\_save\_checkpoint}, che serializza in un file
  binario (compressione LZ4, sidecar SHA3-256) lo \emph{stato completo}
  della run: i metadati di configurazione (checksum della matrice di
  disegno, checksum del vettore $\bm{z}$, coefficiente di ridge,
  tolleranze, budget), il vettore di pesi $\bm{w}$ e la finestra
  $(\bm{s}_i, \bm{y}_i, \rho_i)_{i=1}^{m}$ in ordine cronologico;
\item \textbf{warm start}: i parametri \texttt{initial\_sl},
  \texttt{initial\_yl}, \texttt{initial\_rho} e \texttt{initial\_iter}
  permettono a una nuova chiamata a \texttt{lbfgs} di \emph{riprendere}
  dallo stato checkpointato: la finestra viene pre-popolata (la
  ricorsione a doppio ciclo la trova già pronta all'ingresso nel ciclo
  principale) e la numerazione di $k$ continua da \texttt{initial\_iter}
  (il vincolo \texttt{initial\_iter} $< \texttt{max\_iters}$ è validato
  all'ingresso, \texttt{ValueError} altrimenti).
\end{itemize}

La coerenza del checkpoint con la run corrente è verificata a carico dal
livello \texttt{OPTIMIZE} con tre \emph{gate}: checksum della matrice di
disegno (le posizioni e il feature set non sono cambiati), checksum del
vettore $\bm{z}$ (i target non sono cambiati), coefficiente di ridge
(il funzionale non è cambiato). Se uno dei tre gate fallisce, il
checkpoint è scartato e la run riparte da zero: la finestra $(s,y,\rho)$
di un'altra funzione non è trasferibile. Un checkpoint la cui run è
\emph{già completata} (\texttt{reason} di convergenza, o
$\texttt{iters} \ge \texttt{max\_iters}$) viene usato per caricare
direttamente i pesi finali, saltando l'ottimizzazione intera.

\section{Il contratto di uscita}

A fine run \texttt{lbfgs} restituisce la coppia
$(\bm{x}^{\star}, \texttt{info})$ dove \texttt{info} è un dizionario a
nove chiavi:
\begin{itemize}
\item \texttt{w}: il vettore di pesi ottimizzato (è \emph{lo stesso
  oggetto} di $\bm{x}^{\star}$, non una copia);
\item \texttt{converged} (bool) e \texttt{reason} (una delle tre stringhe
  di \S\ref{sec:lbfgs-restart});
\item \texttt{iters}: il numero assoluto di iterazioni completate
  ($k_{\mathrm{fin}} + 1$);
\item \texttt{f} e \texttt{g\_norm}: il valore del funzionale e la norma
  del gradiente \emph{al punto risultato} (non all'ultima iterazione
  loggata: sono ricalcolati sullo stato finale);
\item \texttt{sl}, \texttt{yl}, \texttt{rho}: la finestra di memoria
  finale — lo stato da ripassare come \texttt{initial\_*} a un eventuale
  warm start successivo.
\end{itemize}

Il livello \texttt{OPTIMIZE} proietta il dizionario sulle cinque chiavi
rilevanti per il contesto (\texttt{hlp\_optimize\_info\_dict}:
\texttt{converged}, \texttt{reason}, \texttt{iters}, \texttt{f},
\texttt{g\_norm}), pubblica \texttt{ctx.w} $= \bm{x}^{\star}$ e il suo
checksum SHA3-256, e salva il checkpoint finale. Il livello
\texttt{GENERATE} che segue consuma $\bm{x}^{\star}$ come vettore di
pesi del modello: è il ponte numerico tra l'ottimizzazione trattata in
questo capitolo e la funzione di valutazione del capitolo 5.

\section{Riepilogo}

L'algoritmo implementato in \texttt{opt\_lbfgs.py} è un L-BFGS con:

\begin{itemize}
\item approssimazione inversa dell'Hessiana a memoria limitata,
  rappresentata dalle $m$ coppie $(\bm{s}_i,\bm{y}_i,\rho_i)$ più
  recenti e applicata al gradiente con la ricorsione a doppio ciclo, con
  il fattore centrale $\gamma\bm{I}$ ricalcolato a ogni iterazione dalla
  coppia più recente (\/\S\ref{sec:lbfgs-twoloop});
\item ricerca del passo a condizioni di Wolfe \emph{forti}, in due fasi —
  espansione (al più 3 iterazioni, passo iniziale $\alpha = 1$,
  $\alpha_{\max} = 10$, fallback a raddoppio) e zoom (al più 20
  iterazioni, fallback a punto medio) — con prossimo punto di prova
  dato dal minimo dell'interpolante cubica di Hermite, costruito sui
  valori \emph{e} sulle derivate prime agli estremi dell'intervallo
  (\/\S\ref{sec:lbfgs-line-search}--\S\ref{sec:lbfgs-hermite});
\item finestra di memoria a rotazione FIFO con condizione di curvatura
  $\bm{y}^{\!\top}\bm{s} > 10^{-15}$ all'inserzione
  (\/\S\ref{sec:lbfgs-history});
\item riavvio a discesa del gradiente quando la direzione non è di
  discesa, e doppi criterio di arresto (stazionarietà del gradiente,
  stall del calo relativo) oltre al budget di iterazioni
  (\/\S\ref{sec:lbfgs-restart});
\item checkpointing periodico e warm start per la ripresa da crash
  (\/\S\ref{sec:lbfgs-checkpoint}).
\end{itemize}

Ogni componente è testata in isolamento dalla suite
\texttt{test\_opt\_lbfgs.py} (24 test): l'identità bit-per-bit della
ricorsione a doppio ciclo, il contratto del dizionario \texttt{info},
la validazione degli input, i tre motivi di arresto, la convergenza su
funzioni sintetiche a difficoltà crescente (quadratica, Rosenbrock,
Zakharov 20D) e la coerenza del warm start (100 iterazioni continue
contro 50 $+$ 50 con salvataggio/ripresa, entro $10^{-4}$).
